\section{Integral length bounds from rank-three Frobenius traces}
\label{sec:rank-three-length}
\hypertarget{source-ordinary-rank-three-length}{}

The descent has identified the ordinary constituents in the two packets. We now return
to the original coefficient ring to measure the length they contribute.
The graded deformation directions give a bound on the
integral lengths in terms of valuations of Frobenius trace
differences in the two descended packets. The bounds enter the odd-prime cohomology window of
Appendix~\ref{sec:cohomology}. The localization kernel maps to the torsion derivations in \eqref{eq:local-length-1}. The final bound \eqref{eq:local-length-11} uses the valuation on the original coefficient field and one representative of each two-sided orbit.

Conjugation by \(D(t)=\operatorname{diag}(t,1,t^{-1})\) gives the grading directly on matrix units. With \((w_1,w_2,w_3)=(1,0,-1)\),
\[
 D(t)E_{ij}D(t)^{-1}=t^{w_i-w_j}E_{ij}.
\] An entry in row \(i\)
and column \(j\), numbered from zero, has weight \(j-i\).
The semidirect product describing the polarization acts on its
Lie algebra through \(X\mapsto-JX^{\mathsf T}J\) on the
quadratic coset, with \(J\) the reversal matrix. This
involution preserves the entry weights. Write \(r_0\) for the
split representation and \(g_*=g_1\) for the mark. The
matrix of \(r(c)\) is fixed. The normalization at the mark
diagonalizes
\[
              r(g_*)-r(cg_*^{-1}c^{-1})
              \quad\text{on the quadratic kernel.}
\]
The difference at \(r_0\) has
diagonal entries \(\Psi(g_*)-\Psi(g_*)^{-1},0,
\Psi(g_*)^{-1}-\Psi(g_*)\). These are distinct in the residue
field by the separator and oddness of \(\ell\).

\subsection{The integral cotangent comparison}

Which part of the integral cotangent module can record the desired length? The grading
isolates it. Let \(O\) be the coefficient DVR, with fraction field \(E\),
uniformizer \(\pi\) and \(v_\pi(\pi)=1\). If \(x:A\to O\) is
the origin of a graded ambient algebra, put \(J_x=\ker x\).
The lattice at weight \(-1\) is the image \(C_{-1}\) of
\(A_{-1}\) in \(J_x/\operatorname{sat}_O(J_x^2)\). The
equations \(F_i\) of weight \(-1\) give a submodule
\(\langle F_i\rangle\) of this lattice.

\begin{construction}
\[
\begin{array}{ccc}
\text{Smith summand}&\text{admissible image of its generator}&\text{length}\\ \hline
O/(\pi^d),\ d<s&\pi^{s-d}O/(\pi^s)&d\\
O/(\pi^d),\ d\ge s&O/(\pi^s)&s\\ \hline
&&\min(d,s)
\end{array}
\]
\caption{One Smith summand in the conormal bound, with \(s=v_\pi(b)\). A point derivation with values in \(O/(b)\) can see at most the smaller exponent.}\label{fig:cotangent-smith-factors}
\end{construction}

\begin{lemma}[Torsion derivations and the saturated conormal]
\label{lem:torsion-derivations-conormal}
Suppose an additive
injection from an \(O\)-module \(M\) sends a class \(m\) to
an \(O/(b)\)-valued point derivation, for \(0\ne b\in O\).
Suppose each derivation kills the \(F_i\) and is the
reduction of an \(O\)-valued point derivation of weight
\(-1\) on \(A\). If the quotient
\(C_{-1}/\langle F_i\rangle\) has finite length, then
\[
 \operatorname{length}_O M\le
 \operatorname{length}_O(C_{-1}/\langle F_i\rangle).
\]
\end{lemma}

\begin{proof}
Every integral point derivation at \(x\)
kills \(J_x^2\) by the product rule and its saturation
by torsion freeness of \(O\). It follows that the injection
factors through
\[
             \operatorname{Hom}_O(C_{-1}/\langle F_i\rangle,O/(b)).
             \tag{1}\label{eq:local-length-1}
\]
Indeed the scalars and every weight other than \(-1\)
already have zero derivative. If the domain lattice
quotient has finite length, the summand \(O/(\pi^d)\)
in its Smith decomposition gives a summand of length
\(\min(d,v_\pi(b))\) in \eqref{eq:local-length-1}. Counting both finite additive
groups in powers of the residue field cardinality gives
the desired length inequality even for an additive
injection. In fact if \((d_i)_i\) are the Smith exponents,
the target has length
\(\sum_i\min\{d_i,v_\pi(b)\}\le\sum_i d_i\).
\end{proof}

\begin{auxlemma}[Annihilation in large multidegree]\label{lem:rees-multidegree-annihilation}
Let \(R\) be a Noetherian \(\mathbb N^s\)-graded algebra generated over
\(R_{\boldsymbol0}\) by the pieces \(R_{\mathbf e_i}\) in the unit
directions. Put \(R_+=\prod_{i=1}^s (R_{\mathbf e_i})\).
Let \(H\subset R\) be a homogeneous ideal and \(J\subset R_{\boldsymbol0}\) an ideal.
If each prime containing \(H\) but not containing \(R_+\) contains \(J\),
there is an integer \(h\ge1\) such that
\[
 J^h(R/H)_{\mathbf d}=0\qquad(d_i\ge h\text{ for every }i).
\]
\end{auxlemma}

\begin{proof}
Every prime over \(H\) contains \(J R_+\), so \(J R_+\subseteq\sqrt H\).
By finite generation \((J R_+)^h\subseteq H\) for some \(h\).
A monomial of multidegree \(\mathbf d\) with every \(d_i\ge h\)
contains at least \(h\) factors from each unit direction. Consequently
\[
 R_{\mathbf d}\subseteq R_+^h,\qquad
 J^h R_{\mathbf d}\subseteq (J R_+)^h\subseteq H.
\]
\end{proof}

\begin{proposition}[Graded cotangent length bound]\label{prop:rank-three-conormal}
Let \(B=A_0\) be complete Noetherian local
and let \(A\) be generated over \(B\) by finitely many
elements of weights \(-2,-1,1,2\). Write \(I\) for the
ideal of the representation chart in \(A\). The ambient
point has a smooth formal chart at \(x\) over \(E\),
with the first row giving the tangent dimensions and
the second row the number of equations of \(I\).
The equations vanish at the section \(x\).
\[
\begin{array}{c|ccccc}
 j&-2&-1&0&1&2\\ \hline
 \text{ambient dimensions}&s&r&m&r'+1&s'+1\\
 \text{equations}&s&r+1&m&r'&s'
\end{array}
\tag{2}\label{eq:local-length-2}
\]
The equations in nonpositive weights span the rational
cotangent spaces of their weights.

Assume that \(B/(I\cap B)\)
is finite over \(O\). Over an algebraic closure of \(E\)
a point is two-sided if some negative and some positive
weight coordinate are nonzero. Over a fixed point other
than \(x\), at least one of the two strictly one-sided
sets with the same restriction to \(B\) must be empty.
At a two-sided point all tangent derivations of \(A/I\)
are generated by the torus tangent. A function \(a\in B\)
vanishes at \(x\) and is nonzero at every two-sided point.

For \(r\ge1\), these hypotheses give
\[
 \operatorname{length}_O(C_{-1}/\langle F_i\rangle)
       \le 2\sum_{\xi\in T}v_\pi\bigl(\xi(a)\bigr),
 \tag{3}\label{eq:local-length-3}
\]
where \(T\) contains exactly one point from each two-sided
torus orbit.
\end{proposition}

\begin{proof}
\proofstep{Slices transverse to the grading torus.}
A point of an orbit can be rescaled to the
slice \(f=1\) for a weight \(-2\) element \(f\) nonzero
on it. Use also a weight \(2\) element \(g\) nonzero on
it. Such even-weight elements are available by squaring.
Only finitely many \(f,g\) from generating sets are
needed. Inverting \(fg\) makes both \(f\) and \(g\) units.

On each slice all even-weight values follow from the
invariant values. Products of pairs of odd-weight values
follow in the same way. The only remaining ambiguity is the action of \(\mu_2\) on the slice. Thus for the tuple of these coordinates it is a common sign:
\[
 (y_i)_i\longmapsto(\varepsilon y_i)_i,\qquad\varepsilon\in\{1,-1\}.
\] The slices are finite over the invariant fibers
and reduced by the tangent hypothesis. This proves the
finiteness of \(T\). It also proves that each generic
invariant point on this locus contributes at most one
copy to any given weight in the unsquared equation quotient.

\proofstep{Squaring the equation list.}
For completeness, start the length calculation by replacing
the \(F_i\) with a Smith list. In the formal chart the linear terms of these equations are independent in their respective coordinates, and the extra relation starts in higher degree:
\[
 \operatorname{lin}(F_i)=\pi^{h_i}y_i\ (1\le i\le r),\qquad \operatorname{lin}(F_*)=0 .
\] The length on the left of
\eqref{eq:local-length-3} equals \(\sum_i h_i\). Over \(E\), eliminate the
weight-zero variables by the \(m\) independent equations
of weight zero. Denote the resulting formal chart map
by \(\phi\). If necessary, replace \(a\) by
\(a+F_i u\) for \(u\in A_1\) with nonzero image under
\(\phi\), so that \(\phi(a)\) is nonzero. This change
keeps every value in \eqref{eq:local-length-3}.

Use the three Rees directions over \(B\) given by
\[
        A_{-2},\qquad A_2,\qquad BF_*^2+A_{-2}^2A_2.
        \tag{4}\label{eq:local-length-4}
\]
In the power series chart write the variables as \(x_i\)
with weights \(w_i\). A finite set \(M_j\) of componentwise
minimal exponent vectors generates the solutions of
\(\sum_i w_i\alpha_i=j\) by addition of weight-zero exponents.
Indeed the minimal set is finite by Dickson's lemma, and any
nonminimal vector above \(\beta\in M_j\) has weight-zero difference
\(\alpha-\beta\). Assign each term of a series to one such \(\beta\).
This gives a finite sum \(\sum_{\beta\in M_j}x^\beta f_\beta\)
with each \(f_\beta\) of weight zero.

Similarly at weight \(2n\)
one can remove a weight-two submonomial for \(n\ge1\): use a
positive weight-two factor when one is present, or two factors
of weight one otherwise. Repeat on the remaining monomial.
Reverse all signs at weight \(-2n\). Grouping the terms shows
\[
                       A_{\pm2n}=A_{\pm2}^{\,n}.
\] Quotient the Rees ring by the
homogeneous ideal of coefficients killed by \(\phi\), and
write \(R\) for the quotient. Choose sections \(\mathbf t\)
consisting of \(F_i^2\) and the weight \(-2\) equations
in the first direction, the squares of the weight \(1\)
equations and the weight \(2\) equations in the second,
and \(F_*^2\) in the third. Regularity of \(a\) on this quotient means that multiplication by \(a\) gives the exact sequence
\[
 0\longrightarrow R\xrightarrow{\ \cdot a\ }R\longrightarrow R/aR\longrightarrow0 .
\]

\proofstep{Preserving the relevant leading coefficients.}
In large multidegree \(d\), the pieces of \(R/(\mathbf t)\)
are finite \(O\)-modules. To check the support, work away
from the ideal generated by products of sections of the
three directions. Choose a section outside the prime
in each direction. Any relation of \(I\cap B\) can
be written as a sum of homogeneous equation multiples.
For even-weight equations, multiply by enough of the
chosen sections to put each term in a Rees degree in
which the equation is a factor.

For an odd-weight equation, pass to its square in the even comparison and then use the implication valid in a prime ideal \(\mathfrak p\):
\[
 z^2\in\mathfrak p\quad\Longrightarrow\quad z\in\mathfrak p .
\]
A section with coefficient \(F_*^2\) in the first
direction can be exchanged with the section carrying
the same coefficient in the third direction after
multiplication by the chosen third section. Thus, with \(J=I\cap B\) and \(R_+\) the product of the
three ideals generated by unit-direction sections, we have
\[
 (\mathbf t)\subseteq\mathfrak p,\quad R_+\not\subseteq\mathfrak p
     \quad\Longrightarrow\quad J\subseteq\mathfrak p\cap B .
\]
This contraction holds on the integral Rees ring.

Suppose now that a characteristic-zero contracted prime
also kills \(A_{-2}A_2\). Nonvanishing of the first and second
sections implies that their coefficients survive outside
the nilradical in the corresponding fiber of \(A/I\).
Indeed, if a power of either coefficient belongs to the
fiber ideal, clearing denominators and making the same
section comparisons forces the power of the section
into the prime. This gives both a negative and a positive
point above the same invariant values. Both points are one-sided. Indeed, at a point over a field the vanishing of all products forces at least one side to vanish:
\[
 A_{-2}(z)A_2(z)=0\quad\Longrightarrow\quad A_{-2}(z)=0\ \text{or}\ A_2(z)=0 .
\] Scaling to the fixed point
keeps their invariant values.

For the fiber of \(x\),
the rational spanning in weight \(-1\) kills all weight
\(-1\) values on a negative one-sided point. In the
saturated square at weight \(-1\) every contribution
has a positive factor or a weight-zero factor vanishing
at \(x\). In weight \(-2\) every additional product has two factors of weight \(-1\). If \(z\) is the point under consideration, then for each such pair
\[
 (ab)(z)=a(z)b(z)=0\qquad(a,b\text{ of weight }-1).
\]
Spanning then also kills the weight \(-2\) values. Thus
the negative point cannot occur at \(x\). At every other
fixed point the one-sided hypothesis gives the exclusion.

\proofstep{Saturation in the three Rees directions.}
Apply Lemma~\ref{lem:rees-multidegree-annihilation} over the
coefficient DVR to \(H=(\mathbf t)\), using the integral prime
contraction just proved. For a fixed integer \(h\ge1\) this gives
\[
      J^h(R/(\mathbf t))_d=0\qquad(d_1,d_2,d_3\ge h).
\] Each nilpotent thickening of
the finite invariant quotient is also finite over \(O\),
by the filtration by powers of its defining ideal.
Adding the equation \(a=0\) leaves only special-fiber support in these
pieces by its assumed nonvanishing at the two-sided
points. These support tests apply to the homology of the
corresponding Koszul lists, which is killed by the list
of sections.

In large multidegrees write \(\chi\) for the alternating
sum of lengths of Koszul homology. The comparison is
\[
   \chi(\mathbf t;R/aR)_d
          \le \chi(a;(R/(\mathbf t))_d).
          \tag{5}\label{eq:local-length-5}
\]
Here \(a\) is regular on \(R\), so the left side equals
\(\chi(\mathbf t,a;R)_d\). Successively taking
kernel and cokernel of multiplication by each section
gives the inequality. The omitted terms have nonnegative
Euler characteristics in large degrees. This positivity,
also used for monotonicity under quotients, follows by
homogeneous prime filtration and induction on the list:
on a prime factor a nonvanishing first section is regular.

If it vanishes, the contribution is the difference of the
two shifted Euler characteristics for the tail. If the replacing section has multidegree \(e\) and is regular on the prime factor \(M\), the calculation uses the sequence
\[
 0\longrightarrow M(-e)\xrightarrow{\ \cdot s\ } M\longrightarrow M/sM\longrightarrow 0 .
\]
Thus only its degree enters the difference of shifted Euler characteristics. If the whole section degree vanishes on the
prime factor, the factor has zero pieces in large degrees
in that direction. For a degree-zero section one can
take the replacement to be \(1\).

\proofstep{The lower bound from the Smith exponents.}
On the left of \eqref{eq:local-length-5}, take leading coefficients in the chart \(\phi\) in the respective negative and positive weighted degrees
\[
 X_d=2d_0+4d_2,\qquad Y_d=2d_1+2d_2 .
\]
Each element in degree \(d\) has at least these totals:
\(F_*\) has weight \(-1\) and no linear term, so its
negative total is at least two. The leading coefficients
form a finite free \(O\)-module acted on by \(B\)
through \(x\). They span the full bidegree over \(E\).
To see the spanning, use representatives of the
individual chart axes in their weights. Group the
weight \(1\) axes in pairs and the weight \(2\) axes
singly on each side. Each third-direction factor takes
two negative blocks and one positive block. Within the
fixed bidegree the products of axis representatives
have independent leading terms by the filtration by
ordinary total degree. The resulting action factors through the vanishing of \(a\) at the point, namely
\[
 R\twoheadrightarrow R/aR\qquad(a(x)=0),
\]

On this quotient \(F_i^2\) acts by \(\pi^{2h_i} y_i^2\).
Splitting off its scalar by Koszul additivity contributes
\(2h_i\) times the Euler characteristic modulo \(\pi\)
of the list with this section omitted. The degree
pieces are torsion-free, so the successive
\(\pi\)-quotients count the same monomials. Positivity
also applies to the list with all scalars removed.
The Hilbert series for the two leading degrees is
\[
 \frac{1}{(1-X)^r(1-X^2)^s(1-Y)^{r'+1}(1-Y^2)^{s'+1}}.
\]

The remaining sections multiply it by
\((1-X^2)^{r-1+s}(1-Y^2)^{r'+s'}(1-X^4Y^2)\).
After cancellation the first numerator factors are
\((1+X)^{r-1}(1+Y)^{r'}\). Separate the two degrees in the numerator by adding and subtracting \(Y^2\):
\[
 1-X^4Y^2=(1-Y^2)+Y^2(1-X^4).
\]
The second summand has bounded \(X\)-degree after
cancellation. The first summand has eventual coefficient
\(2^{r+r'-1}\). Thus the left side of \eqref{eq:local-length-5} is at least
\(2^{r+r'}\sum_i h_i\).

\proofstep{The upper bound from two-sided orbits.}
On the right of \eqref{eq:local-length-5} the kernel and cokernel lengths on
the \(O\)-torsion cancel. On the free part the Euler
characteristic is the valuation of the determinant of
multiplication by \(a\). It is the sum over the geometric
invariant points with the multiplicities in the generic
degree piece. At these points multiplication by a negative and by a positive even-weight unit permits shifts between the corresponding weight pieces: if \(u\) is such a unit of weight \(w\),
\[
 M_j\xrightarrow[\sim]{\ \cdot u\ } M_{j+w},\qquad m'\longmapsto u^{-1}m'\ \text{for the inverse}.
\] Each module in \eqref{eq:local-length-4}
then becomes its full weight piece.

The generic Rees
piece is a quotient of the corresponding weight piece
with ordinary equations \(F_i^2,F_*^2,E_i^2\), the
weight \(2,-2\) equations, and the weight-zero
equations, where \(E_i\) are the weight \(1\) equations.
Replacing an equation \(q\) by \(q^2\) at most doubles
each multiplicity bound, by the two pieces of the
filtration by \(q\), with one grade shifted. The
unsquared orbit bound above gives multiplicity at
most \(2^{r+1+r'}\) after the replacements. The factor
of two in \eqref{eq:local-length-3} follows on comparing with the lower bound.
\label{cor:ordinary-rank-three-length-1}

\end{proof}

\subsection{The split point and the localization kernel}

To apply the cotangent comparison we must identify the part of the localization kernel
it detects. Begin at the split point with rational coefficients, then pass to finite
torsion depth.

\begin{lemma}[Rational vanishing at the normalized split point]
\label{lem:rank-three-split-vanishing}
For the split representation at the Gaussian fixed-inertia chart of
Proposition~\ref{prop:gaussian-rank-three-chart} or the general
crystalline chart of Proposition~\ref{prop:rank-three-crystalline-local-base}, the normalized rational point derivations satisfy 
\[
 \operatorname{Der}_{x,j}^{\rm norm}=0\qquad(j=-2,-1,0),
\]
 where the subscript denotes the weight at the split point \(x\). In these
weights the equations therefore span the rational conormal
spaces in the hypotheses of Proposition~\ref{prop:rank-three-conormal}.
\end{lemma}

\begin{proof}
Write \(\Gamma\) for the host Galois group, and
\(\Psi=\epsilon_\ell^{\,e}\psi\). On the quadratic kernel
the weight \(j\) adjoint space of \(r_0\) has scalar action
\(\Psi^j\). On \(\Gamma\) its eigencharacters are
\(\Psi^j\) and \(\eta\Psi^j\), on the \(+1\) and \(-1\)
eigenspaces, respectively, of \(-JX^{\mathsf T}J\).
The involution sends \(E_{ij}\) to \(-E_{2-j,2-i}\) for \(0\le i,j\le2\). A two-element orbit contributes one eigenvector of each sign. A fixed matrix position contributes one negative eigenvector. Counting the orbits in each weight gives
\[
 \begin{array}{c|rrrrr}
 j&-2&-1&0&1&2\\ \hline
 \dim (\mathrm{Ad}\,r_0)_j&1&2&3&2&1\\
 \dim (\mathrm{Ad}\,r_0)_j^{+}&0&1&1&1&0\\
 \dim (\mathrm{Ad}\,r_0)_j^{-}&1&1&2&1&1
 \end{array}
\]

The rational nonpositive-weight groups used at the split
point vanish by the following rank-one argument. At
weight zero restrict to the quadratic kernel, so that
both characters are trivial. The local condition gives
an unramified additive character everywhere. For the
Gaussian chart this follows from vanishing on the
residual-kernel inertia and finite index there, and
from the unramified weight-zero condition at \(\ell\).
For the crystalline chart, inertia at \(\ell\) lies in
the quadratic kernel and the finite local conditions on
the trivial character are unramified. Class field theory
makes the resulting homomorphism zero, since the class
group is finite. Restriction to an open subgroup is
injective on rational cohomology by corestriction.

\proofstep{The negative-weight strict local condition.}
For negative \(j\), write \(j=-n\) with \(n>0\). The corresponding characters have the form
\[
 \epsilon_\ell^{ej}\theta=\epsilon_\ell^{-en}\theta,\qquad\theta\text{ finite}.
\]
At \(\ell\) the Gaussian local condition kills the class;
a crystalline tangent at the split point has finite local
condition, which is zero for a negative Tate twist of a
finite character. Indeed its de Rham quotient by
\(\mathrm{Fil}^0\) and its invariant space are zero.
Restriction to a finite field killing the finite local
type gives the same finite condition by averaging over
the finite quotient. Away from \(\ell\) an unramified
class with these coefficients vanishes because Frobenius
minus one on its inertia invariants is invertible. Even the full local rational cohomology vanishes here.

For this one-dimensional character space \(V\) away from \(\ell\), the local Euler characteristic formula and duality give
\[
 \dim H^1_v(V)=\dim H^0_v(V)+\dim H^0_v(V^\vee(1))=0 .
\] Thus the
global class becomes everywhere locally split on a finite
abelian extension \(L\) killing \(\theta\) and containing
\(\mu_\ell\). The field is totally complex since
\(\ell\) is odd. Take \(S\) over \(L\) to contain the
places above \(\ell\) and the ramification of the cocycle
and coefficients. The classes under consideration are
unramified outside such a finite \(S\).

\proofstep{Rational vanishing from the height and duality inputs.}
Use the finiteness part of the height estimate of
Section~\ref{sec:cohom-height} at the fixed field \(L\)
and exponent \(en+1\). For a \(G_{L,S}\)-module \(M\) write \(\operatorname{Sha}^i(M)\)
for the kernel of localization in degree \(i\), with all
places of \(S\) included. Poitou--Tate gives
\[
 \operatorname{Hom}\!\left(\operatorname{Sha}^1(\Q_\ell/\Z_\ell(-en)),
                                     \mathbb Q/\mathbb Z\right)
       \simeq \operatorname{Sha}^2(\Z_\ell(en+1)).
\]
The height input makes the right side finite, so the strict
torsion kernel on the left has finite order.
Extra places of \(S\) preserve the required finiteness bound
on this kernel: in the degree-two localization sequence
the new terms are local terms, where the class in the
kernel has zero image.

Equivalently an everywhere locally
split cocycle in degree one factors through the group
unramified outside the places above \(\ell\).
Finite coefficient extension preserves this finiteness.
A nonzero rational class would give torsion classes of
unbounded order: multiply a rational class by
arbitrary inverse powers of the uniformizer and project
to the divisible coefficient module. Let \(V\) be the rational coefficient representation and \(T\) its lattice. The sequence of coefficients identifies the kernel of projection on cohomology:
\[
 \ker\bigl(H^1(V)\to H^1(V/T)\bigr)=\operatorname{im}\bigl(H^1(T)\to H^1(V)\bigr).
\]
The image is a lattice in its rational span by finiteness of the residual global cohomology. Local splitting persists
under projection. This contradicts finiteness of the
strict torsion kernel.

For a cocycle with values in an
algebraic coefficient closure, a finite extension contains
its values. To see this, use the countable family of finite
extensions containing the field of definition of the action.
Their preimages under the cocycle are closed subgroups,
and Baire's theorem gives one with open preimage. For a cocycle \(z\) let \(H\) be the open preimage obtained in this way. If \(g_1,\ldots,g_r\) represent \(G/H\), adjoining the finitely many coordinate values of \(z(g_i)\) suffices by
\[
 z(g_i h)=z(g_i)+g_i z(h)\qquad(h\in H).
\] This proves the nonpositive-weight vanishing
over both fields of coefficients.

An infinitesimal normalized point at \(r_0\) now has
zero derivative in each of these weights. Its matrix
derivative times \(r_0(g)^{-1}\) is a cocycle, by the
product law with square-zero coefficients. A coboundary
is infinitesimal conjugation by \(1+\varepsilon X\).
Fixing \(r(c)\) and the normalization at \(g_*\) imposes the simultaneous conditions
\[
 X=-JX^{\mathsf T}J,\qquad X\text{ diagonal}.
\] Such an
\(X\) is a multiple of \(\operatorname{diag}(1,0,-1)\),
which centralizes the split polarized representation.
Since matrix entries separate infinitesimal deformations,
the chart derivative is zero. Dualizing, the equations
span the rational cotangent spaces in \eqref{eq:local-length-2}. In the
Gaussian power series presentation this can also be
tested over \(O\): a rational vector in the annihilator
would have an integral nonzero multiple annihilating
the equations and would give a normalized point over
\(O[\varepsilon]/(\varepsilon^2)\).
\hfill\mbox{\cite{bloch-kato-finite,milne-duality}}\label{cor:ordinary-rank-three-length-2}

\end{proof}

\paragraph{Transfer at every torsion depth.}
Rational vanishing deals with the generic fibre. A length comparison requires a map
that detects classes even at finite depth.

\begin{proposition}[Integral transfer at all coefficient depths]
\label{prop:rank-three-integral-transfer}
The weight-minus-one localization kernel with coefficients
\((E/O)(\Psi^{-1})\), on either of the two designated hosts,
injects at each torsion depth into the derivations in condition
\textup{\eqref{eq:local-length-1}}. Lifting uses the Gaussian chart or the crystalline
local base, respectively. Let \(G\) be this kernel over the coefficient ring \(O\) with
uniformizer \(\pi\) and absolute ramification index \(e_O\). A uniform
length bound from \eqref{eq:local-length-3} passes to the full kernel:
\[
 (\forall h\ge1,\ \operatorname{length}_O G[\pi^h]\le L)
       \quad\Longrightarrow\quad \operatorname{length}_O G\le L .
\]
The normalized bound is divided by \(e_O\) to use \(v_\ell\).
\end{proposition}

\begin{proof}
\proofstep{The weight-minus-one cocycle.}
Let \(W=(E/O)(\Psi^{-1})\) and \(b\in O\setminus\{0\}\).
For brevity put \(W[b]=(O/(b))(\Psi^{-1})\), using the
isomorphism \(u\mapsto u/b\bmod O\). A class in the
rank-one localization kernel killed by \(b\) has a
representative valued in \(W[b]\). In fact if \(bz\)
is the coboundary of \(u\in W\), choose a vector \(u'\)
with \(bu'=u\) and subtract its coboundary. The class
in finite-coefficient global cohomology is unique, since
\(\Psi(g_*)^{-1}-1\) is a unit and thus \(W^\Gamma=0\).

On an inertia subgroup on which \(\Psi=1\), the action on the coefficient is trivial. Thus a coboundary on that subgroup satisfies
\[
 (da)(i)=i\cdot a-a=0,
\]
and the finite representative vanishes there. At a decomposition group where some
\(\Psi(g)^2-1\) is a unit, a trivializing vector in \(W\)
for this representative is killed by \(b\). Indeed
its product with \(\Psi(g)^{-1}-1\) is killed by \(b\).
Consequently the local class already vanishes with
coefficients in \(W[b]\).

The weight \(-1\) adjoint embedding of \(W[b]\) uses
\(T=E_{10}-E_{21}\), where \(E_{ij}\) is the elementary
matrix in row \(i\) and column \(j\). Write
\(g_*'=cg_*^{-1}c^{-1}\). Normalize the cocycle \(z\)
by making
\[
    L(z)=\Psi(g_*)z(g_*)-\Psi(g_*')z(g_*')
\]
zero. Subtract the coboundary of
\(L(z)/(\Psi(g_*)^{-1}-\Psi(g_*))\), which is defined
since the denominator is a unit. Also \(z(c)=0\)
because \(c^2=1\), \(\Psi(c)=1\) and \(2\) is a unit.
The representation
\[
             (1+\varepsilon z(g)T)r_0(g),
       \qquad \varepsilon^2=0,\quad b\varepsilon=0,
             \tag{6}\label{eq:local-length-6}
\]
therefore satisfies the two normalization conditions.
Write \(D_{10}(g)\) for the coefficient of \(\varepsilon\) in the \((1,0)\) entry of \eqref{eq:local-length-6}. Since \(T_{10}=1\), the entry calculation gives
\[
    D_{10}(g)=\Psi(g)z(g),\qquad
    z(g)=\Psi(g)^{-1}D_{10}(g)\quad(g\in\Gamma).
\]
Thus the resulting derivatives determine the normalized cocycle, and addition of cocycles adds the derivatives entrywise.

\proofstep{Integral lifting on the Gaussian chart.}
For the Gaussian chart quotient by the torsion-depth
local conditions at \(\ell\) and at the exceptional
tame places where \(\Psi(g)^2-1\) is a unit on some
element of the decomposition group. These are precisely
the places where weight \(-1\) must be a local coboundary
in \eqref{eq:local-length-6}. At the other tame places the representation
kills the residual-kernel inertia, on which
\(\Psi=1\). The normalized representation therefore
defines a point of its universal algebra. In the ambient power series ring centered at the split point, let \(X_i\) be its coordinates.

For each coordinate of weight \(-1\) choose \(d_i\in O\) lifting the given derivative in \(O/(b)\); set \(d_i=0\) in every other weight. The point derivation
\[
  D(f)=\sum_i d_i\,\frac{\partial f}{\partial X_i}(0)
\]
is integral. Modulo \((b)\) it gives the original infinitesimal point, so \(D(F)\equiv0\pmod b\) for every equation \(F\) of the chart. The original
localization group on the absolute Galois group
descends to \(G_{\Q,S}\), with \(S\) the primes dividing
\(\ell N\), because its cocycles vanish on all inertia
outside \(S\). Thus \eqref{eq:local-length-1} bounds every \(b\)-torsion
subgroup by the same cotangent length.

\proofstep{Integral lifting on the crystalline chart.}
For the general-data chart the global condition kills
the residual-kernel inertia at each place away from
\(\ell\). The datum makes \(\Psi\) trivial on those
subgroups. For a class \(z\) in this localization kernel the condition at each prime dividing \(\ell N\) is on the full group:
\[
 \operatorname{res}_{D_p}[z]=0\qquad(p\mid\ell N).
\] Otherwise it is zero already on
inertia. In either case the finite representative vanishes
on the subgroup to be killed, including above primes
ramified only in the quadratic field.

At \(\ell\) the mark's unit separator gives a trivializing vector in \(O/(b)\). Lift a conjugating matrix for this vector to a weight \(-1\) matrix \(Y\) over \(O\), choosing its sign to recover the given coboundary. The integral local lift is
\[
 (1+\varepsilon Y)r_0(g)(1+\varepsilon Y)^{-1}
    =r_0(g)+\varepsilon\bigl(Yr_0(g)-r_0(g)Y\bigr),
       \qquad \varepsilon^2=0
\]
on the positive component, with the polarized action extending it. Being conjugate to the split point, it lies on the crystalline local quotient. The ambient global chart is a relative
power series algebra over that quotient, as described
below. Sending its remaining coordinates to arbitrary
integral lifts of their derivatives in weight \(-1\)
gives an integral point derivation of the ambient algebra.
Its reduction comes from \eqref{eq:local-length-6} and vanishes on every
global equation. This is the lift required by \eqref{eq:local-length-1}.
Lifting on the crystalline local quotient first is
necessary for this use of saturation.

\proofstep{The union over all coefficient depths.}
The divisible module \(W\) is the union of its
finite torsion submodules, and a continuous cocycle
into it has finite image by compactness. Thus the
localization group is a union of its \(\pi^h\)-torsion
subgroups. For the localization group \(G\) the torsion filtration and its common length bound \(L\) satisfy
\[
 G_h=G[\pi^h]\subseteq G_{h+1},\qquad \operatorname{length}(G_h)\le L\quad(h\ge1).
\] Additivity gives
\[
 \operatorname{length}_O(G_{h+1}/G_h)
 =\operatorname{length}_O G_{h+1}-\operatorname{length}_O G_h .
\]
Every strict inclusion contributes at least one; the sum of these increments is at most \(L\). The chain stabilizes and its union \(G\) has length at most \(L\). Dividing by the
ramification index of \(O/\Z_\ell\) takes the bound
in \eqref{eq:local-length-3} from \(v_\pi\) to \(v_\ell\). The procedure
keeps the packet bound independent of torsion depth
and of the residue-field and coefficient degrees.
\label{cor:ordinary-rank-three-length-3}

\end{proof}

\subsection{The Gaussian inertia chart}

The integral comparison is ready to use once the local presentation is known. We
construct it first for the Gaussian point, where the deformation fixes inertia.

\begin{proposition}[The Gaussian fixed-inertia presentation]
\label{prop:gaussian-rank-three-chart}
Under the Gaussian mark and its local conditions, the graded
presentation has the dimensions and equations of \textup{\eqref{eq:local-length-2}},
with a finite degree-zero ring as required for
Proposition~\ref{prop:rank-three-conormal}. Let \(n_{-1}\) be the number of weight-minus-one variables and \(G\) the
localization kernel of Proposition~\ref{prop:rank-three-integral-transfer}.
The coordinate map in \eqref{eq:local-length-6} gives
\[
                         n_{-1}=0\quad\Longrightarrow\quad G=0 .
\]
\end{proposition}

\begin{proof}
\proofstep{The fixed-inertia local problem.}
Use \(S=\{v:v\mid\ell N\}\), the Auxiliary field over \(\mathbb Q(i)\) quotient
\(\Gamma=G_{\Q,S}\), and the one mark from
Section~\ref{sec:marked-hosts}. At \(\ell\) the local
representation preserves a flag with the three ordered
inertia characters \(\Psi,1,\Psi^{-1}\). At an exceptional
tame place it is diagonal up to a change of basis
congruent to the identity, with the prescribed inertia
characters. On the other tame decomposition groups the inertia subgroup killed by the representation is prescribed by
\[
 N_v=\ker(\eta|_{I_v})\cap\ker(P|_{I_v}),\qquad P=\bar\Psi .
\] This is the subgroup on which the
coefficient character is trivial there. In both tame
cases \(I_v/N_v\) has order prime to \(\ell\).
The local determinant or polarization condition is the
one inherited from the semidirect product.

\proofstep{Lifting the local conditions across an extension.}
These local conditions lift across a graded Artinian
small extension. Here the kernel is killed by the
maximal ideal and is supported in one weight. At
\(\ell\), use the eigenspaces of the separator to
choose the flag basis, with the lines labelled by their
reductions. Let \(T\) be the lifted separator and \(\alpha_i\) its labelled eigenvalues. The projector onto its \(i\)-th eigenspace is integral, because
\[
 e_i=\prod_{j\ne i}\frac{T-\alpha_j}{\alpha_i-\alpha_j},\qquad \alpha_i-\alpha_j\in O^\times\quad(i\ne j).
\]
Thus the necessary changes of flag basis with the prescribed entry weights lift. In the flag basis lift the torus
action by taking the given inertia characters and
lifting a Frobenius in the semidirect torus. The
procyclic unramified quotient permits this choice,
also when \(\eta\) acts on the torus. Lift the
upper triangular entries in order of root height.
At each step let \(u(g,h)\) be the multiplication error in
the small-extension kernel of that root weight \(j>0\).

Associativity of the lifted products gives the two-cocycle identity for their additive defect:
\[
 g\,u(h,k)-u(gh,k)+u(g,hk)-u(g,h)=0 .
\] The separator \(x\) has
\(\bar\epsilon_\ell(x)=\eta(x)=1\) and \(P(x)^2\ne1\).
On \(\operatorname{Hom}(V_j,k(1))\) it acts by \(P(x)^{-j}\),
which is nontrivial for both possible positive weights \(1,2\).
Therefore
\[
 H^2(\Q_\ell,V_j)\simeq H^0(\Q_\ell,V_j^\vee(1))^\vee=0.
\]
Write \(u=da\) for a continuous cochain \(a\). Subtracting
\(a\) from the selected root lifts removes the multiplication
error while keeping the other weights fixed. At each exceptional tame place, the separator kills invariants and dual invariants in nonzero weights. Local duality and the Euler characteristic in residue characteristic different from \(\ell\) then give 
\[
 h^2=0,\qquad h^1=h^0+h^2=0,
\]
 where \(h^i\) denotes the local dimension in the selected nonzero weight. A canonical
eigenbasis there diagonalizes the inertial action,
and the decomposition group preserves these lines
with the action prescribed by \(\eta\).

Lifting its
Frobenius in the torus and changing back to the original
basis gives the desired lift. At the remaining tame
places lift the finite prime-to-\(\ell\) action of
\(I_v/N_v\) by averaging. Its lifts are conjugate by
a matrix in the kernel. Let \(r_I\) be the lifted inertia action and \(F\) a first Frobenius lift. Reduction identifies the two inertia maps, since the residual representation already respects conjugation:
\[
 \overline{F r_I(i) F^{-1}}=\overline{r_I(\mathrm{fr}\,i\,\mathrm{fr}^{-1})}.
\] The averaging argument conjugates
one to the other by a kernel matrix \(C\). Replacing \(F\)
by \(CF\) gives
\[
         (CF)r_I(i)(CF)^{-1}
                    =r_I(\mathrm{fr}\,i\,\mathrm{fr}^{-1}).
\]
This is the required Frobenius-inertia relation. The
extension of the finite inertia group by the unramified
procyclic quotient is generated by inertia and one
Frobenius with those relations. All the linear
corrections in the kernel can be taken in the chosen
weight, because the residual adjoint action preserves
the grading.

\proofstep{The tangent and obstruction dimensions.}
We give the dimension computation controlling the global
relations. Fix a weight \(j\), let \(d=\dim V_j\),
let \(d_\infty=\dim V_j^{c=1}\), and write \(g_i,h_{v,i}\)
for the dimensions of the global and local cohomology
groups. For the space of permitted local cocycle classes the inclusion is
\[
 L_v\subseteq H^1(\mathbb Q_v,V_j).
\] At a tame place
\(\dim L_v=h_{v,0}\). Indeed restriction to inertia
of a cocycle vanishing on \(N_v\) is a coboundary by
averaging. Let \(T\) be Frobenius on \(V_j^{I_v}\). Rank-nullity identifies the dimension of its coinvariants with that of its invariants:
\[
 \dim\operatorname{coker}(T-1)=\dim\ker(T-1).
\]

At exceptional places the same
formula holds by the specified inertia and the
vanishing of \(H^1\) in nonzero weights.
At \(\ell\), \(L_\ell=0\) for \(j<0\), is unramified
for \(j=0\), and is all of \(H^1\) for \(j>0\).
The tangent assertion for the flag follows directly
from the ordered characters. A negative-weight
entry comes only from changing the flag basis, a
weight-zero entry can vary its unramified character,
and an upper-weight cocycle gives a flag-preserving
deformation with fixed diagonal.

Let \(t_j\) be the dimension of the Selmer tangent
space and \(o_j\) the dimension of its obstruction
space. If \(\lambda\) is the rank of the localization
map to \(\bigoplus_v H^1(\Q_v,V_j)/L_v\), then
\[
 \begin{split}
 t_j&=g_1-\lambda,\\
 o_j&=\sum_v (h_{v,1}-\dim L_v)-\lambda
        +g_2-\sum_v h_{v,2}.
 \end{split}
 \tag{7}\label{eq:local-length-7}
\]
The first part of \(o_j\) is the degree-one localization
cokernel. The rest is the kernel of localization in
degree two in the cochain complex recording a global
multiplication error and its local trivializations.
The degree-two localization map is surjective by the
Poitou--Tate end sequence.

In fact
\(H^0(\Q,V_j^\vee(1))=0\) for each \(j\): for
nonzero \(j\) use the preceding separator, and for
\(j=0\) use an element of inertia at \(\ell\) with
nontrivial residual cyclotomic value. The global
Euler characteristic is
\(g_1+ d_\infty=g_0+g_2+d\).
The local Euler characteristic is
\(h_{v,1}=h_{v,0}+h_{v,2}+d\,1_{v=\ell}\).
Subtracting the two lines of \eqref{eq:local-length-7} gives
\[
           t_j-o_j=g_0-d_\infty+\dim L_\ell-h_{\ell,0}.
           \tag{8}\label{eq:local-length-8}
\]
For \(j\ne0\) the global and \(\ell\)-local invariants
are zero by the separator. In positive weights
\(h_{\ell,1}=d\) because \(h_{\ell,2}=0\).
At weight zero \(g_0=d_\infty=1\).
Consequently the five defects, in increasing weight,
are \(0,-1,0,1,1\).

For clarity, these are bounds for a presentation of
the normalized deformation problem. The residual
representation is fixed, the flags in the local
conditions have the marked residual labels, and the
inertia conditions are invariant under a change of
frame. At the global origin any infinitesimal
cocycle can have its value at \(c\) and the
off-diagonal derivative of the marked difference removed
by a coboundary. If \(z\) is the adjoint cocycle and
\(dY(g)=g\cdot Y-Y\), then
\[
 z(c)+c\cdot z(c)=0,\qquad d\bigl(-z(c)/2\bigr)(c)=z(c).
\]
Subtract this coboundary first. The remaining normalization
uses a \(c\)-invariant matrix \(Y\).

In the eigenbasis of the
marked operator with eigenvalues \(\lambda_i\), let \(D=\operatorname{diag}(\lambda_i)\). To realize the required off-diagonal change \(Z\) in the derivative solve
\[
 [D,Y]_{ij}=(\lambda_i-\lambda_j)Y_{ij}=Z_{ij},\qquad
 Y_{ij}=\frac{Z_{ij}}{\lambda_i-\lambda_j}\quad(i\ne j).
\]
The denominators are units at the mark. The prescribed change belongs to the invariant off-diagonal space, on which this commutator is invertible. The remaining conjugators centralize the split
point. Thus normalized tangents contain
one representative of every Selmer class.

The cohomology of every open subgroup of \(G_{\Q,S}\)
with finite \(\ell\)-torsion coefficients is finite in
degree one. This is the finite-dimensional tangent
condition for pro-representability of the framed
deformation problem. The frame supplies a complete
Noetherian universal ring. Its matrix entries,
after subtracting their split values, give
homogeneous variables. Choose homogeneous entries whose classes form a basis of the relative cotangent space modulo the maximal ideal. With \(r\) such entries \(x_i\), the resulting presentation over the coefficient ring \(O\) has the map
\[
 O[[X_1,\ldots,X_r]]\twoheadrightarrow R,\qquad X_i\longmapsto x_i .
\] That choice generates by complete
Nakayama. The variables have weights in
\([-2,2]\). The ideal of relations is homogeneous,
as follows by applying the torus character with an
indeterminate unit and comparing weights in each
finite quotient.

\proofstep{The relation count for the completed global presentation.}
Here is the relation bound underlying this presentation.
For a graded small extension, choose a set-theoretic
lift of the matrices and local homomorphism lifts.
Let \(I\) denote the square-zero extension kernel, \(c\) the global multiplication error and \(b_v\) the local correction at \(v\). In the mapping-cone convention of \eqref{eq:local-length-7}, associativity and the local lifting equations give
\[
  dc=0,\qquad \operatorname{res}_v(c)-db_v=0 .
\]
If the global lift changes by \(\alpha\), let \(\lambda_v\) be the difference between the two permitted local lifts, read as a local tangent cocycle. Its class lies in the prescribed local subspace \(L_v\). The change of the error pair is
\[
   \bigl(d\alpha,(\operatorname{res}_v\alpha-\lambda_v)_v\bigr).
\]
A local change of basis gives the special case \(\lambda_v=d\beta_v\).

In the relative complex with the specified local tangent conditions, all these changes give the same obstruction class, as reflected in the quotient by \(L_v\) in \eqref{eq:local-length-7}. If a functional \(u\) on the extension kernel kills the obstruction class \(o\), its pushout has obstruction
\[
 o_{\rm pushout}=u_*(o)=0 .
\]
Hence there is a global lift with the required local conditions. The normalization at \(c\)
and the mark can be restored there by a kernel
conjugation as above. Apply this to the dual of the
space of minimal relations of one weight in the
power series presentation. The calculation can be
made in a finite quotient inducing the same space
of minimal relations, by Artin--Rees. A functional
giving a lift would let the universal ring map back
to the corresponding nonzero extension of its
Artinian quotient.

If the two maps out of the power series ring differ by \(\delta_i\in I\) on the variables, then for each minimal relation \(F\) their difference is
\[
 F(\mathbf X+\boldsymbol\delta)-F(\mathbf X)
     =\sum_i (\partial_i F)(\mathbf X)\,\delta_i=0.
\]
The Taylor expansion ends after its linear part because \(I^2=0\). For the minimal relations each linear coefficient lies in the maximal ideal, so 
\[
 (\partial_i F)\delta_i\in\mathfrak m I=0\qquad\text{for every }i.
\]
 The section would thus kill the selected relation even in the extension where its class is nonzero, a contradiction. Hence the dual
of the minimal relation space injects into
the obstruction space. This proves the bound of
\(o_j\) relations in weight \(j\).

If \(n_j\) is the number of chosen variables, then
\(n_j\ge t_j\). By \eqref{eq:local-length-8} the relation bounds can be
padded by zero equations to
\[
 (n_{-2},\ n_{-1}+1,\ n_0,\ n_1-1,\ n_2-1).
\]
There is at least one variable in each positive
weight by \eqref{eq:local-length-8}. In case \(n_{-1}=0\) every
localized class has zero image in the coordinate
derivations in \eqref{eq:local-length-6}, so the desired length is zero.
Otherwise these variables and equations have
exactly \eqref{eq:local-length-2}. Take \(A\) to consist of finite
sums of homogeneous series in the ambient power
series ring. Its split formal chart has the
specified weight dimensions. Homogeneous projection
shows the ideal of the quotient is generated in
\(A\) by the homogeneous equations. The invariant
ring in the quotient is finite by the Gaussian
fixed-mark argument, and it also makes the
rescaling to integral points possible. This gives
the chart to which \eqref{eq:local-length-3} applies.
\hfill\mbox{\cite{milne-duality}}\label{cor:ordinary-rank-three-length-4}

\end{proof}

\subsection{The crystalline chart for general data}

For the general point the local calculation at the coefficient prime uses a
crystalline chart. In the rank-three trace estimate the order of \(\psi\) is prime to \(\ell\),
so its values can be taken over an unramified
coefficient DVR. Normalized length is invariant
under further coefficient extension. Work over
the general datum with \(\ell\nmid D_K h_K\),
where \(h_K\) is the class number. Away from \(\ell\)
kill \(I_v\cap\ker\eta\cap\ker P\); the split lift
kills it because the finite inertia character has
order prime to \(\ell\). We may work with the
finite set of places above \(\ell N|D_K|\), using
the normal closure of the killed subgroups in the
global deformation problem.

\paragraph{The base at \(\ell\).}
To keep the codimension count integral we must control saturation at the coefficient
prime. The following elementary calculation is the local ingredient.

\begin{auxlemma}[A coefficient test for saturation]\label{lem:crystalline-coefficient-test}
Let \((B,\mathfrak n)\) be a local ring with \(\mathfrak n^h=0\).
Suppose \(P,y\in B[t]/(t^n)\) satisfy \(Py=0\) and the image of \(P\)
in \((B/\mathfrak n)[t]/(t^n)\) is \(t^b\) times a unit, where
\(0\le b<n\). Evaluation at the origin then satisfies
\[
                       n>hb\quad\Longrightarrow\quad y(0)=0\quad\text{in }B .
\]
\end{auxlemma}
\begin{proof}
Write \(y=\sum_{i<n}y_i t^i\). We prove
\[
  i<n-vb\quad\Longrightarrow\quad y_i\in\mathfrak n^v
  \qquad(0\le v\le h).
\]
For \(v=0\) the assertion holds in \(B\). Suppose it holds at \(v<h\)
and truncate modulo \(t^{n-vb}\). The truncated coefficients of \(y\)
lie in \(\mathfrak n^v\), so \(P\) acts on their classes in
\(\mathfrak n^v/\mathfrak n^{v+1}\) by its residue \(t^b u\), with \(u\)
invertible. Cancelling \(u\) in the truncated polynomial module gives
\[
 t^b\sum_{i<n-vb}[y_i]t^i=0\pmod {t^{n-vb}}.
\]
Consequently \([y_i]=0\) for \(i<n-(v+1)b\), as required.
At \(v=h\), \(0<n-hb\) puts \(y_0\) in \(\mathfrak n^h=0\).
\end{proof}

\begin{construction}
\[
\begin{array}{ccc}
 B &\longleftarrow& B[t]/(t^n)\\[-2pt]
 \downarrow\ {\scriptstyle\bmod\,\mathfrak n}&&
    \downarrow\ {\scriptstyle\bmod\,\mathfrak n}\\
 E&\longleftarrow&E[t]/(t^n)\\
 \multicolumn{3}{c}{\text{horizontal arrows evaluate at }t=0}\\[2pt]
 \multicolumn{3}{c}{\ker\bigl(B[t]/(t^n)\longrightarrow
       B\times_E E[t]/(t^n)\bigr)=t\mathfrak n[t]/(t^n)}
\end{array}
\]
\caption{The test algebra joining the given point to the unramified curve. Evaluation and reduction have common target in the residue field. Formal smoothness lifts the compatible pair across the resulting nilpotent kernel.}\label{fig:crystalline-thickening-test}
\end{construction}

\begin{proposition}[The saturated crystalline local chart]
\label{prop:rank-three-crystalline-local-base}
For the general-data host with \(\ell\nmid D_K h_K\) and
\(\psi\) of order prime to \(\ell\), kill
\(I_v\cap\ker\eta\cap\ker P\) on inertia away from \(\ell\).
Use the prime through
the split point in the radical of the ordered-inertia equations
saturated by \(\ell D_0D_1\). Its formal characteristic-zero chart is crystalline, including at
the split point. The dimensions \(d_j\) of the homogeneous coordinate
spaces in \eqref{eq:rank-three-crystalline-chart} are
\[
                (d_{-2},d_{-1},d_0,d_1,d_2)=(1,2,3,4,2).
\] On the selected global chart the test-curve
comparison gives a finite degree-zero algebra, and the
two-sided characteristic-zero orbits have integral representatives.
\end{proposition}

\begin{proof}
\proofstep{The ordered-inertia equations and their saturation.}
Let \(U_\ell\) be the unrestricted framed
local ring and \(H=\ker\eta\cap\ker\tau\) in
the local group, where \(\tau\) is the finite
character in \(\Psi|_{I_\ell}=\epsilon_\ell^e\tau\).
For \(s\in H\cap I_\ell\) put \(q_i(s)=\Psi(s)^{1-i}\),
\(0\le i\le2\), and \(S_i(s)=r(s)-q_i(s)1\).
Impose the characteristic polynomial with roots
\(q_0(s),q_1(s),q_2(s)\), and impose
\[
          S_0(s_0)S_1(s_1)S_2(s_2)=0
          \qquad(s_i\in H\cap I_\ell).
\]

Choose \(s\in H\cap I_\ell\) of infinite
cyclotomic image. Its three \(q_i(s)\) are
distinct. With \(c_i=\prod_{j\ne i}(q_i(s)-q_j(s))\),
the spectral projector onto the \(i\)-th
eigenline is \(c_i^{-1}\prod_{j\ne i}S_j(s)\)
over the fraction field. Take \(\mathrm{fr}\in H\)
generating the unramified quotient and let
\(T_i\) be the trace of
\(\prod_{j\ne i}S_j(s)r(\mathrm{fr})\).
The two adjacent differences are
\[
       D_i=c_{i+1}T_i-
            \epsilon_\ell(\mathrm{fr})c_iT_{i+1}
       \qquad(i=0,1).
\]
The local ideal is the prime through the split
point of the radical of the ordinary equations
saturated by \(\Delta=\ell D_0D_1\).
We explain the use of saturation in both directions.

\proofstep{The ordered crystalline equations after inverting the adjacent factors.}
In characteristic zero the ordinary equations give
an ordered flag. For the first line and last
quotient, use a variable inertia element in
the first and last factors of the triple product
respectively and the chosen \(s\) in the other two.
The projectors isolate the first line and the last quotient of the ordered flag. On those one-dimensional graded pieces the operator thus has the two scalar restrictions
\[
 q_0\,\mathrm{id}\quad\text{on the first line},\qquad q_2\,\mathrm{id}\quad\text{on the last quotient}.
\] The middle value follows from the trace.
These flags are unique with the ordered characters
because the eigenvalues at \(s\) are distinct.
The decomposition group preserves the flag with
the conjugate-dual action on the quadratic coset.
On \(H\), write \(\nu_1,\nu_2,\nu_3\) for the unramified twists. The ordered graded characters restrict to
\[
 \nu_1\epsilon_\ell^e,\qquad\nu_2,\qquad\nu_3\epsilon_\ell^{-e}.
\]
In the ordered basis \(T_i/c_i\) is the \(i\)-th
Frobenius value.

Fix a finite local characteristic-zero coefficient
algebra \(B\) on the saturated quotient. One can
factor its point through a finite local algebra
\(B'\) of that quotient such that any element of
\(B'\) killed by \(\Delta\) maps to zero in \(B\).
Indeed in the generic local ring of the saturated
quotient, \(\Delta\) is a nonzerodivisor. Apply
Artin--Rees to its principal submodule and the
powers of the maximal ideal \(\mathfrak m\).
Writing \(R\) for this local ring, choose \(c\) from Artin--Rees,
so \(\mathfrak m^h\cap\Delta R\subseteq\mathfrak m^{h-c}\Delta R\)
for \(h\ge c\). Since \(\Delta\) is a nonzerodivisor,
\[
 \Delta b\in\mathfrak m^h\quad\Longrightarrow\quad
 b\in\mathfrak m^{h-c}.
\]
Choose \(h\) large enough that this last ideal maps to zero in \(B\). Taking
\(B'\) to be this \(\mathfrak m^h\)-quotient
gives the assertion. Its residue field is
finite over the coefficient fraction field,
since the point has values in a finite
algebra.

\proofstep{The vanishing required by the saturation.}
On \(H\), a ratio of adjacent characters is an
unramified twist of \(\epsilon_\ell^e\).
For a positive Tate exponent \(a\), the
quotient of \(H^1\) by its finite subspace
is dual to the finite \(H^1\) of the Tate
dual, by local duality. For \(a>1\) the Tate dual has exponent \(1-a<0\). Its finite local group vanishes, so local duality gives
\[
 H^1/H_f^1=0\qquad(a>1).
\] For \(a=1\), write the coefficient as
\(\epsilon_\ell\otimes\theta\) with \(\theta\) unramified.
The Tate dual is \(\theta^\vee\). Its finite
\(H^1\) is the Frobenius coinvariant quotient.
Dualizing that quotient makes the possible nonfinite
part the kernel of \(\theta(\mathrm{fr})-1\).
In terms of the adjacent characters this scalar is
\[
 \theta(\mathrm{fr})-1
  = \frac{(\kappa_i/\kappa_j)(\mathrm{fr})
                          -\epsilon_\ell(\mathrm{fr})}
          {\epsilon_\ell(\mathrm{fr})}.
\]
The denominator is a unit.

Over \(B'\) a class in this
kernel is killed by the indicated factor of \(\Delta\),
hence is killed by \(\Delta\), so its image in \(B\)
is zero by the choice of \(B'\). This gives finite
extensions on the two adjacent subquotients.
For the outer character ratio the positive exponent is at least two, so the preceding local dimension calculation gives
\[
 2e\ge2\quad\Longrightarrow\quad H^1/H_f^1=0
\]
for that ratio.
There is no additional obstruction for the
full upper triangular representation. Once
the upper two-dimensional block is crystalline,
tensor it by the inverse of the last character.
Its two successive characters have positive
exponents. The finite \(H^1\) maps onto the
finite \(H^1\) of its last quotient by exactness
and the dimension formula for the de Rham
quotient. The remaining outer-line class is
finite as just proved. Thus the full point
over \(B\) is crystalline on \(H\).

\proofstep{Recovering the crystalline quotient, including the split point.}
Conversely take a crystalline deformation on \(H\)
over a local Artin coefficient algebra reducing
to the split lift. The Frobenius after a full
residue-field period has three distinct slopes
at the split lift and hence three separate
eigenlines over the Artin algebra. Order the three Frobenius eigenlines \(D^{(1)},D^{(2)},D^{(3)}\) by increasing slope. The initial subspaces for which crystallinity must be checked are
\[
 D_1=D^{(1)},\qquad D_2=D^{(1)}\oplus D^{(2)} .
\] To verify this with
nilpotent coefficients, regard the filtered
Frobenius module as a vector space over the
coefficient field. Its multiset of Hodge
weights is the split multiset repeated
according to the length of the Artin
algebra. The Newton number of each initial
eigenline sum is the sum of the corresponding
lowest Hodge weights. The induced filtration on a subspace of this dimension has
at least the sum of the smallest Hodge weights, counted with
those multiplicities.

If \(D_i\) is either initial eigenline
sum, and \(h_{\min}(D_i)\) is that minimum, then
\[
 h_{\min}(D_i)\le t_H(D_i)\le t_N(D_i)=h_{\min}(D_i).
\]
The middle inequality is weak admissibility of the full
object on the Frobenius-stable \(D_i\). Equality, together
with the inherited inequalities on all its subobjects, makes
\(D_i\) weakly admissible. The crystalline equivalence
gives the desired flag on the representation.
The graded terms are rank one over the
Artin algebra, with fixed Tate exponents
and unramified factors on \(H\). Freeness
can be checked by reduction and the
underlying vector space dimension.

On the overlap of two base changes the two pulled-back ordered flags have the same defining property. Uniqueness therefore gives the descent compatibility, with \(p_1,p_2\) the two projections:
\[
 p_1^*\mathcal F=p_2^*\mathcal F .
\] Its finite inertial
characters lift those at the split point,
because the order of a finite inertia
quotient is invertible over a characteristic-zero
coefficient field. This proves the ordinary
equations on every such deformation.

For a single crystalline point as above,
\(\Delta\) need not be a unit. To see that
it kills the saturated ideal, let \(\Delta^k a\)
lie in the ordinary ideal and use the following
deformation of the test algebra. Take the local algebra \(B\) with residue coefficient field \(E\) and maximal ideal \(\mathfrak n\) of nilpotence index at most \(h\):
\[
 B/\mathfrak n=E,\qquad \mathfrak n^h=0 .
\] Choose \(n>2kh\). On
\(C=E[t]/t^n\), vary the split lift by
an unramified torus character with value
\(\operatorname{diag}(1+t,1,(1+t)^{-1})\)
on a generator of the unramified quotient.
On \(\mathrm{fr}\) the power of \(1+t\)
has nonzero exponent. Each adjacent difference either has nonzero constant term or differs from a nonzero multiple of \(t\) by higher terms.

For the latter case in characteristic zero use
\[
 (1+t)^a-1\equiv at\pmod{t^2},\qquad a\ne0 .
\]
Its order in \(t\) is therefore one. The image of \(\Delta\)
is \(t^j\) times a unit with \(j\le2\).
The local crystalline functor at the split
lift is formally smooth. In the filtered
Frobenius description one lifts the invertible
Frobenius and the filtration by direct summands
across a nilpotent extension. Across the strict successive layers the Hodge and Newton numbers are additive. Writing \(D_i\) for the successive quotients, the equality starts with
\[
 t_H(D)=\sum_i t_H(D_i)=\sum_i t_N(D_i)=t_N(D).
\]
Strictness and weak admissibility of the layers give the subobject inequalities as well. Equivalence with
crystalline representations recovers the action
and one lifts the frame on its free coefficient
module. The fixed finite type on descent
is unchanged across the nilpotent extension.
This is the local smoothness of the crystalline
quotient theorem.

On the unramified
quadratic kernel the conjugate-dual descent
datum is retained by order-two averaging
on the characteristic-zero lifting and
change-of-basis spaces. The given \(B\)-point and the unramified twist over
\(C\) reduce to the same split lift. They therefore give
a point over \(B\times_E C\). The coefficient map
\[
       B[t]/t^n\longrightarrow B\times_E C,\qquad
             f(t)\longmapsto (f(0),\bar f(t))
\]
is surjective with kernel \(t\mathfrak n[t]/(t^n)\).
The kernel is nilpotent. Formal smoothness with the
polarized descent retained gives a lift across this map.

\proofstep{Keeping the integral local equation after saturation.}
Let \(P\) be the image of \(\Delta^k\)
and \(y\) the image of \(a\) in this
thickening. The ordinary equations give
\(Py=0\) and \(P\bmod\mathfrak n\) is
\(t^{jk}\) times a unit. At the two possible positive weights the nilpotence bound stays strictly below the truncation order:
\[
 j\le2\quad\Longrightarrow\quad hjk\le2kh<n .
\] Lemma~\ref{lem:crystalline-coefficient-test} with \(b=jk\) gives \(y(0)=0\). Evaluation at \(t=0\) recovers the given \(B\)-point, so that point kills \(a\), as required for every element of the saturated ideal.
The saturation consequently gives the
crystalline formal germ at the split lift,
including tangent points over dual numbers.
The completed germ is smooth and therefore
has a single component through that point.
Its closure defines the integral local base.

Conjugation congruent to the identity preserves
the split component. Indeed the action by the
universal congruence matrix on a domain
component lands in a formal power series domain
over that component. The preimage of its zero
prime under the conjugation map is contained in
the original component prime by specialization
to the identity. If \(P\) denotes the original component prime and \(Q\) the contracted zero prime, this gives a chain of primes
\[
 Q\subseteq P,\qquad P\text{ minimal}\quad\Longrightarrow\quad Q=P .
\] In particular the conjugate of
the split representation over the integral
dual numbers factors through the chosen local
quotient, as used in the transfer.

\proofstep{The local weight directions on the selected component.}
With characteristic-zero adjoint coefficients, the dimension at the split point is
\[
 \dim H^1_f(V_j)=h^0_j+\dim\bigl(D_{\rm dR}(V_j)/\mathrm{Fil}^0\bigr).
\]
Changing the actual frame contributes
\(\dim V_j-h^0_j\). The de Rham quotient has
dimension \(\dim V_j\) in positive weights and
zero in other weights. This follows from the
Tate twists \(\epsilon_\ell^{ej}\) on the
weight spaces, with only finite character
factors. The local framed smooth chart therefore
has dimensions
\[
                    (1,2,3,4,2)
                    \tag{9}\label{eq:rank-three-crystalline-chart}
\]
in increasing weight. Homogeneous matrix
entries give compatible axes. In the formal
completion over the fraction field, their
differentials span because a tangent annihilating
all entries gives the split deformation.

Choose
a homogeneous basis among them and change formal
coordinates by the inverse function theorem.
This gives the graded local chart, with the
same weights on the ambient integral
presentation. The integral component injects into
the smooth completed chart, since its defining
prime is the contraction of the prime of that
chart. Matrix entries map to homogeneous functions.
Thus for a relation with homogeneous components \(F_j\), projection to any fixed weight in the completed chart yields
\[
 \sum_j F_j=0\quad\Longrightarrow\quad F_i=0\quad\text{in weight }i .
\] By injectivity on the
integral component, the component ideal is closed
under these weight projections. The integral
relative presentation inherits this completed
grading.
\hfill\mbox{\cite{bloch-kato-finite,colmez-fontaine,kisin-crystalline}}\label{cor:ordinary-rank-three-length-5}

\proofstep{Inertia on test curves.}
The ordinary flag at the generic point of
the same component has the full prescribed
inertia characters. On \(H\cap I_\ell\)
these are given by the ordinary equations.
For a character ratio \(\delta\) trivial on \(H\cap I_\ell\), the remaining inertia datum factors through the finite quotient
\[
 I_\ell\twoheadrightarrow I_\ell/(H\cap I_\ell)\xrightarrow{\ \delta\ }\operatorname{im}(\delta).
\]
Specialization at the split lift fixes this finite factor.
More explicitly the eigenspaces of \(r(s)\)
at the separating element \(s\) label the
flag across any specialization in which
the \(q_i(s)\) stay distinct, including
the split lift over the fraction field.
The ratios at other inertia elements
have finite order and hence are constant
on the component on this labelled
chart once their value at the split
point is fixed. The ordered flag with
its prescribed action is a closed
condition on the projective flag
scheme over the integral local base.

Let \(Z\) denote that closed flag locus and \(S\) its integral base. Properness makes the image closed, so
\[
 \eta_S\in\operatorname{im}(Z\to S)\quad\Longrightarrow\quad\operatorname{im}(Z\to S)=S .
\]
Each field specialization thus has a flag after extending its residue field.
This also follows by extending the
generic flag over valuation rings.
For a map to a complete local DVR, write \(A\) for the
matrix of the separator. Its distinct reduced eigenvalues
lift by Hensel's lemma to roots \(a_i\) with unit differences.
The projectors
\[
                   e_i=\prod_{j\ne i}
                        \frac{A-a_j}{a_i-a_j}
\]
give the integral eigenlines. The ordered flag after field
extension is the sum of these lines in the prescribed residual
order, so it descends to the original valuation ring. This supplies the marked
ordinary condition even when the DVR
kills the coefficient uniformizer.

\proofstep{Applying the residual polynomial comparison to this chart.}
Now apply the finite-end statement of
Section~\ref{sec:marked-hosts} to
each such characteristic-\(\ell\) DVR
point of the crystalline global ring.
Away from \(\ell\), the inertia elements to be killed form the intersection
\[
 \ker(\eta|_{I_v})\cap\ker(\Psi|_{I_v}) .
\]
For the general datum this agrees with the prescribed residual kernel. The statement
makes the characteristic polynomial
split on a fixed subgroup of finite
index. The finite-index argument in
that section makes it split on the
whole quadratic kernel, and the mark
then gives triangularity with the
split diagonal. Thus every diagonal
coordinate minus its split value, and
every balanced monomial with a
positive and a negative entry, vanishes
on every test curve.

As on the auxiliary field over \(\mathbb Q(i)\), the curve test puts each such generator \(z\) in the radical of the uniformizer ideal. Equivalently, in the indicated local ring \(R\),
\[
 z\in\sqrt{\varpi R}\quad\Longleftrightarrow\quad z^m\in\varpi R\ \text{for some }m\ge1 .
\] Complete Nakayama gives
finiteness of the degree-zero algebra
over \(O\) on the selected chart.
The same valuation interval as in
the Gaussian case gives an integral
representative of every two-sided
characteristic-zero orbit. Here the
fact that this point is on the
crystalline quotient supplies the
finite condition at \(\ell\) for
its infinitesimal deformations.
\label{cor:ordinary-rank-three-length-6}

\end{proof}

\paragraph{The relative global presentation.}
With the integral base in hand we can count the remaining global variables and
relations relative to it.

\begin{lemma}[The normalized relative presentation]
\label{lem:relative-rank-three-presentation}
Over the integral local base of Proposition~\ref{prop:rank-three-crystalline-local-base},
the completed relative presentation after normalization at
\(c,g_*\) has the counts of \eqref{eq:local-length-10}. In increasing order of
weights \(-2,-1,0,1,2\), its numbers of coordinates and equations are
\[
 \begin{aligned}
 (n_j)_j&=(a_{-2}+1,a_{-1}+2,a_0+3,a_1+4,a_2+2),\\
 (r_j)_j&=(a_{-2}+1,a_{-1}+3,a_0+3,a_1+3,a_2+1).
 \end{aligned}
\]
These counts, its rational chart and the integral derivative lift
apply to the same quotient.
\end{lemma}

\begin{proof}
\proofstep{Separating the local and relative global presentations.}
Before normalizing at \(c,g_*\), the relative
deformation problem fixes a framed representation
at \(\ell\) on the chosen local quotient. In the
global presentation it records the local
conjugacy comparing the two frames. For a
residual weight space \(V_j\), write
\(g_i,h_{v,i}\) as in \eqref{eq:local-length-7}, and let \(\lambda\)
be the rank of localization into the full
\(H^1\) at \(\ell\) and the quotient by the
unramified conditions at the other places.
Those away conditions have dimension \(h_{v,0}\)
by averaging over the residual inertia quotient.
At \(\ell\) a global relative tangent has
zero cohomology class. Its local trivializing
vector can still vary by local invariants.

After taking the global changes of frame
into account, the relative tangent dimension is
\[
                a_j=g_1-\lambda+h_{\ell,0}-g_0 .
\]
Here the cancellation of the common change of
frame can also be seen by setting the local
conjugacy equal to the identity at \(\ell\)
and allowing the global frame to vary. For
the obstruction space the relative complex gives
\[
 h_{\ell,1}+\sum_{v\ne\ell}(h_{v,1}-h_{v,0})
                  -\lambda+g_2-\sum_v h_{v,2}.
\]
To justify the last sum, the general datum
gives an element with \(\eta=-1\),
\(P=1\) and \(\bar\epsilon_\ell=1\), from
ramification of \(K\) away from \(\ell N\).
Multiplying it by \(c\) gives an element
acting trivially on all \(V_j\) but with
cyclotomic value \(-1\). On the Tate-twisted dual the element just constructed acts as minus the identity. Hence in characteristic zero
\[
 v\in H^0(\mathbb Q,V_j^\vee(1))\quad\Longrightarrow\quad v=-v\quad\Longrightarrow\quad v=0 .
\]
Poitou--Tate then gives the claimed surjectivity in degree two.
The local and global Euler formulas now give
obstruction dimension \(a_j+\dim V_j^{c=1}\).

\proofstep{The relative equations and their weights.}
The relative ring is a quotient of a power
series algebra over the local crystalline
quotient with \(a_j\) additional variables
of weight \(j\). Indeed the relative tangent
coordinates lift a basis modulo the
local base and its maximal ideal. The
small-extension computation for the Gaussian
chart applies with the lift at \(\ell\)
prescribed by this base. One records its
change of frame in degree one of the
relative complex. The tame lift still
uses the finite inertia quotient.
For each weight \(j\), write \(\mathrm{Rel}_j\) for the space of relative minimal relations and \(\mathrm{Obs}_j\) for the corresponding relative obstruction space. The small-extension test gives
\[
 \mathrm{Rel}_j^\vee\hookrightarrow\mathrm{Obs}_j,\qquad\dim\mathrm{Rel}_j\le\dim\mathrm{Obs}_j .
\] This proves the
claimed relation bound over the local base.
Represent the additional coordinates by
matrix entries, so their weights stay in
\([-2,2]\). Zero equations pad each weight
to the bound.

\proofstep{The eight normalization equations.}
There are eight normalization equations.
Applied at \(c\), the relation \(c^2=1\) imposes the following symmetry on the matrix \(Y\) of the quadratic coset:
\[
 Y=JY^{\mathsf T}J .
\] The entries in positions
\((0,0),(1,1),(0,1),(1,0),(0,2),(2,0)\)
therefore suffice to set \(Y=1\).
Once this is done the marked difference is
\(Z-JZ^{\mathsf T}J\) with \(Z=r(g_*)\).
The off-diagonal equations reduce to the
\((0,1)\) and \((1,0)\) entries, one of
each weight \(1,-1\). Thus the normalization
adds \((1,2,2,2,1)\) equations in
increasing weight. The combined chart has
\[
\begin{array}{c|ccccc}
j&-2&-1&0&1&2\\ \hline
\text{ambient dimensions}&
a_{-2}+1&a_{-1}+2&a_0+3&a_1+4&a_2+2\\
\text{equations}&
a_{-2}+1&a_{-1}+3&a_0+3&a_1+3&a_2+1 .
\end{array}
\tag{10}\label{eq:local-length-10}
\]

We retain the entire integral local base
when forming the relative power series
ring. Its formal chart over the fraction
field is obtained by adjoining the same
relative variables to \eqref{eq:rank-three-crystalline-chart}. Passing to the
finite sums of homogeneous components
commutes with the quotient by the
homogeneous equations by projection on
each weight. Thus the smoothness,
equation counts and integral derivative
lifting required for \eqref{eq:local-length-3} and \eqref{eq:local-length-1} refer
to the same chart.
\hfill\mbox{\cite{milne-duality}}\label{cor:ordinary-rank-three-length-7}

\end{proof}

We now have the local and global models used in
Proposition~\ref{prop:rank-three-conormal}. The Gaussian
model and the crystalline model are used in separate
applications. In each one the bound is read on the original
coefficient ring.

\subsection{Frobenius trace bounds for the original Galois representations}

Return to the representations before passage to the marked field. The comparison on
either chart now expresses the length bound in terms of traces in the corresponding
packet.

\begin{proposition}[A length bound in terms of Frobenius trace differences]
\label{prop:rank-three-original-trace-bound}
Use the packet over the original field and the normalized Gaussian
or crystalline chart of Propositions~\ref{prop:gaussian-rank-three-chart}
and~\ref{prop:rank-three-crystalline-local-base}, respectively, with the
relative presentation of Lemma~\ref{lem:relative-rank-three-presentation}.
The normalized
length \(s_\ell(\psi,e)\) satisfies \textup{\eqref{eq:local-length-11}}: its bound is
twice the sum of valuations of the selected Frobenius trace
differences over representatives of the two-sided orbits.
For one character kernel \(J_0\), its rank-three contribution \(S(J_0)\) at the selected test prime satisfies 
\[
 S(J_0)\le 2v_\ell(A_{J_0}(p)),\qquad p\ne\ell.
\]

The Gaussian large-conductor application uses its separate
inert-prime packet. The packet exponents are those of the
original hosts.
\end{proposition}

\begin{proof}
We check the remaining point conditions of \eqref{eq:local-length-3}.
On a fixed point, the restriction to the
quadratic kernel is diagonal. Each of its
three characters differs from the corresponding
split character by a character unramified at
every finite place of \(K\) and congruent
to one. Here are the inertia checks for
that statement. At a tame place of residual
kernel type the diagonal character factors
on inertia through a group of order prime
to \(\ell\). Let \(m\) be the order of the finite inertia quotient. Reduction distinguishes the character values because \((m,\ell)=1\), so
\[
 \zeta,\zeta'\in\mu_m,\quad\bar\zeta=\bar\zeta'\quad\Longrightarrow\quad\zeta=\zeta'.
\] At a tame place of prescribed
inertia type the local eigenbasis has its
residual labels. The same argument using
the local eigenbasis at \(\ell\) gives the
ordered labels there.

In the general-data chart, on the inertia subgroup where the ordered polynomials are imposed, let \(\chi_i\) denote the three diagonal characters and \(\chi_i^0\) their prescribed values. Take a separating element \(h\) and label the three distinct roots \(\lambda_i=\chi_i(h)=\chi_i^0(h)\) in the fixed frame. Work over the characteristic-zero field of this fixed point. For any \(g\) in that subgroup, equality of the traces prescribed at \(h^a g\) gives
\[
 \sum_{i=1}^3 \lambda_i^a\bigl(\chi_i(g)-\chi_i^0(g)\bigr)=0
              \qquad (a=0,1,2).
\]
The determinant of these three equations is \(\prod_{i<j}(\lambda_j-\lambda_i)\ne0\). Thus each labeled value agrees, even when \(g\) itself has coinciding roots.

By unramified class field theory the quotient character \(\xi\)
factors through the class group of \(K\). Its image is finite and
congruent to one, and hence has \(\ell\)-power order. Thus
\[
 |\operatorname{im}\xi|\mid h_K,\qquad
 |\operatorname{im}\xi|=\ell^u,\qquad \ell\nmid h_K
 \quad\Longrightarrow\quad \xi=1.
\] For the Gaussian
case \(h_K=1\). The fixed value at complex
conjugation determines the representation
on the other coset. Hence the only
fixed point is \(x\). This suffices for
the one-sided condition away from \(x\).

\proofstep{The tangent criterion for the two-sided points.}
At a two-sided point of the Gaussian chart,
the block exclusion in Section~\ref{sec:host-descent}
gives an irreducible restriction to \(K\).
For general data, the two possibilities for the restriction are
\[
 V\text{ irreducible of dimension }3,\qquad\text{or}\qquad V=W_{\rm outer}\oplus U_{\rm middle}
\]
with an irreducible outer plane and a middle line. In both charts
restrict the marked integral representative
at \(\ell\) to obtain the finite adjoint
local condition. For the Gaussian chart
this was the fixed-inertia derivative
calculation of Section~\ref{sec:host-descent}.

For the general datum it follows on the
crystalline local quotient by the
crystalline extension criterion, after
the same finite base change that
trivializes the local type. Away from
\(\ell\), the inertia condition gives
the finite local condition by restriction
and averaging. Thus the adjoint class
of any tangent to a two-sided point
vanishes by the rational Selmer results
of that section. The constraints from complex conjugation and the marked difference leave only the following line of coboundary matrices:
\[
 X=u\operatorname{diag}(1,0,-1).
\]
It produces exactly the torus tangent.

\proofstep{Continuity and the completed integral chart.}
These point and derivative arguments apply
on the graded algebra. A degree-zero
point is integral on the finite invariant
quotient. Nonzero products of opposite
weights give the strict interval of
allowable torus rescalings, so the full
matrix coordinates can be taken in the
maximal ideal of a finite coefficient
DVR away from their split values. A first
order lift of this point has the same
extension to the complete chart. For
each generator of nonzero weight its value
is the rescaled integral value plus a
nilpotent term.

In a dual-number ring, a product of \(n\) rescaled
values \(a_i+\delta b_i\) has the expansion
\[
 \prod_{i=1}^n(a_i+\delta b_i)
   =\prod_{i=1}^n a_i+
       \delta\sum_{i=1}^n b_i\!\prod_{j\ne i}a_j .
\]
The \(a_i\) are in the maximal ideal, and the finitely
many \(b_i\) have one fixed lower valuation bound. Both
coefficients therefore tend to zero uniformly in such products
as \(n\) increases. At the
finite invariant quotient, a point
derivation on its local coefficient
algebra is continuous by this same
product rule on powers of the maximal
ideal. Hence completing the chart
gives a normalized dual-number
deformation. Its matrix entries still
separate the derivations. This is the
chart version of the torus tangent
condition in \eqref{eq:local-length-3}.

\proofstep{Applying the conormal bound to the selected trace.}
Take \(a\) to be the trace at the
selected Frobenius minus the split
trace. It has weight zero and
vanishes at \(x\). The packet bounds
of Section~\ref{sec:host-descent}
show it is nonzero at each of
the two-sided points. Apply \eqref{eq:local-length-3}
and the integral transfer \eqref{eq:local-length-1} at
every torsion depth. With \(T\)
one representative of each
two-sided orbit, the normalized
rank-one length satisfies
\[
 s_\ell(\psi,e)
 \ \le\ 2\sum_{\xi\in T}
   v_\ell\!\left(\operatorname{tr}r_\xi(\mathrm{Frob}^{-1})
                  -\operatorname{tr}r_0(\mathrm{Frob}^{-1})\right).
 \tag{11}\label{eq:local-length-11}
\]
An intertwiner on the quadratic
kernel is diagonal in the marked
bases. For an irreducible triple,
polarization then makes it a scalar
times a grading-torus element.
For the plane plus line, the
middle diagonal entry can be
changed to the square root of the
product of the two outer entries.
This preserves the intertwiner.
So in either case the injection
from orbits to the packet indices
in \eqref{eq:local-odd-d} applies to precisely the
representatives in \eqref{eq:local-length-11}.

\proofstep{Normalized valuations on the original packets.}
For the rank-three trace estimate sum \eqref{eq:local-length-11} over the
characters with one kernel. At a test prime \(p\ne\ell\) the common congruence integer bounds the sum over characters with this kernel by
\[
 2v_\ell(A_{J_0}(p)).
\]
Two fixed distinct test primes
give the bound at all admissible
\(\ell\), as in the arithmetic sum
following \eqref{eq:local-odd-l}. The cardinality
and valuation factors are at the
original levels \(N\) and \(N|D_K|\).
Once \(K\) and the pair of test
primes are fixed, the exponent
depends on neither \(N,e\) nor the
residual prime or coefficient
ring.

For the case over the Gaussian field, the fixed
Gaussian field and the inert
prime \(p\le 12\ell N\) give
\(s_\ell(\psi,1)\log\ell\ll
(\ell N)^C\) at the large
dyadic conductor, using only the
rank-three packet. The smaller
conductor goes through the
height estimate as in the
odd-prime count. These deductions
use the height and marked
support inputs of the previous
sections. The enlargement of the
marked field enters neither
packet exponent.
\label{cor:ordinary-rank-three-length-8}
\end{proof}
